% 倍长中线 引入：△ABC，D 是 BC 中点，AB=5, AC=7. 延长 AD 到 E，DE=AD，连 CE。
\documentclass[tikz,border=4pt]{standalone}
\usepackage{ctex}
\usepackage{amsmath}
\usetikzlibrary{calc, decorations.markings}
\begin{document}
\begin{tikzpicture}[scale=1.0, thick,
  tick/.style={postaction={decorate, decoration={markings,
    mark=at position 0.5 with {\draw (-1.5pt,-3pt) -- (1.5pt,3pt);}}}},
  ticktwo/.style={postaction={decorate, decoration={markings,
    mark=at position 0.5 with {\draw (-2.5pt,-3pt) -- (-0.5pt,3pt);
                               \draw (0.5pt,-3pt)  -- (2.5pt,3pt);}}}}]
  % AB=5, AC=7, BC=any (e.g.6). 用余弦定理: cos B = (25+36-49)/(2*5*6)= 12/60=0.2
  % A 在 B 上方. B=(0,0), C=(6,0). A: AB=5 angle B=acos(0.2)≈78.46°
  \coordinate (B) at (0,0);
  \coordinate (C) at (6,0);
  \coordinate (A) at ({5*cos(78.46)}, {5*sin(78.46)}); % ≈(1, 4.9)
  \coordinate (D) at (3,0);
  % E: 延长 AD, DE=AD. E = 2D - A
  \coordinate (E) at ($2*(D) - (A)$);
  \draw (A)--(B)--(C)--cycle;
  \draw[red, dashed] (A)--(D);
  \draw[red, dashed] (D)--(E);
  \draw[blue, dashed] (C)--(E);
  % BD, DC tick
  \draw[tick] (B)--(D);
  \draw[tick] (D)--(C);
  % AD, DE ticktwo
  \node[left] at ($(A)!0.5!(B)$) {\footnotesize $5$};
  \node[right] at ($(A)!0.5!(C)$) {\footnotesize $7$};
  \node[above] at (A) {$A$};
  \node[below left] at (B) {$B$};
  \node[below right] at (C) {$C$};
  \node[below] at (D) {$D$};
  \node[below right] at (E) {$E$};
\end{tikzpicture}
\end{document}
